\chapter{Databricks SQL and SQL Warehouses}
\index{Databricks SQL}\index{SQL warehouse}\index{serverless SQL}\index{Photon}%
\index{materialized view}\index{result cache}\index{query profile}\index{system tables}%
\begin{objectives}
\begin{itemize}
    \item Explain SQL warehouse architecture: serverless, pro, and classic types, and how Photon accelerates queries.
    \item Write warehouse-grade SQL against Unity Catalog tables, including time travel and window functions.
    \item Choose between materialized views, streaming tables, and scheduled Gold tables for derived data.
    \item Size warehouses with min/max clusters, auto-stop, and channel selection; read the query profile.
    \item Attribute cost with \texttt{system.billing.usage} and warehouse monitoring.
    \item Design a warehouse topology that isolates BI bursts from engineering backfills.
\end{itemize}
\end{objectives}

\begin{crosscloud}
Databricks SQL is the platform's \textbf{warehouse-serving layer}; each cloud fields a serverless
SQL engine with the same promises and different mechanics.

\vspace{2mm}
{\footnotesize
\begin{tabular}{@{}p{3.0cm}p{3.4cm}p{3.2cm}p{3.2cm}@{}}
\toprule
\textbf{Capability} & \textbf{Databricks} & \textbf{AWS} & \textbf{Azure / GCP / Fabric} \\
\midrule
SQL compute & SQL warehouse (serverless/pro/classic) & Redshift Serverless / Athena & Synapse serverless / BigQuery / Fabric Warehouse \\
Engine & Photon (vectorized) & Redshift engine / Trino & Synapse / Dremel / Fabric engine \\
Autoscale & Multi-cluster load balancing & Concurrency scaling / per-query & Autoscale / slots / capacity smoothing \\
Idle behavior & Auto-stop, serverless scale-to-zero & Pause / per-query billing & Auto-pause / per-slot / capacity \\
Derived tables & Materialized views, streaming tables & Redshift MVs / Athena CTAS & Synapse MVs / BigQuery MVs \\
Cost visibility & \texttt{system.billing.usage} & CUR + CloudWatch & Cost Management / billing export / Capacity Metrics \\
\bottomrule
\end{tabular}}

\vspace{1mm}
\textbf{Snowflake} is the closest conceptual peer---virtual warehouses with per-second billing and
result caching---but Databricks SQL queries the \emph{same} Delta tables your pipelines write,
with no copy into a proprietary store.
\end{crosscloud}

\section{Week 6 Roadmap and Mental Model}
The first five chapters built and optimized tables; Week 6 serves them. Databricks SQL (DBSQL) is the warehouse surface: SQL editors, SQL warehouses, query federation, and the serving layer beneath AI/BI (Chapter 7). The engineering questions are operational: which warehouse type, what size, how does it scale, what does it cost, and how do you know a query is slow for a fixable reason \cite{databricksSQL,databricksSQLWarehouses}.

The mental model: \textbf{a SQL warehouse is a load-balanced fleet of query clusters over your Delta tables}. You submit SQL; the warehouse compiles it (Photon), pulls data from cloud storage through the cache, scales clusters between its min/max bounds under concurrency, and records everything in system tables. There is no data movement into the warehouse---the storage is the lake you already govern.

\begin{definitionbox}
\textbf{A SQL warehouse} is SQL-optimized compute for Databricks SQL: a managed set of clusters with a size (2X-Small to 4X-Large), a min/max cluster range, an auto-stop timer, and a type (serverless, pro, classic) controlling startup speed, features, and price.
\end{definitionbox}

\begin{keyterms}
\textbf{Photon}: the vectorized execution engine. \textbf{Serverless warehouse}: instant-start, fast-autoscaling compute in Databricks' account. \textbf{Query profile}: the per-operator execution breakdown with rows, bytes, and time. \textbf{Result cache}: reuse of a query's results when query and data are unchanged. \textbf{Materialized view}: a managed, incrementally refreshed precomputed result.
\end{keyterms}

\section{Monday: Warehouse Types and Photon}
\subsection{Serverless, Pro, Classic}
\textbf{Serverless} warehouses start in seconds, scale aggressively, run in Databricks' cloud account, and support the newest features (including predictive optimization integration and AI functions). \textbf{Pro} adds classic-cloud control with fast start. \textbf{Classic} trades startup latency for the most configuration control and lowest DBU rate. Default to serverless for BI; drop to pro/classic when networking or custom initialization requires the data plane.

\subsection{Photon Under the Hood}
Photon re-implements Spark's execution hot path in vectorized native code: columnar batches, SIMD instructions, cache-aware operators \cite{databricksPhoton}. The practical consequences: scans, joins, and aggregations over Delta are dramatically faster than vanilla Spark SQL, and the speedup grows with data scanned---which is why Chapter 5's layout work multiplies this chapter's engine.

\begin{notebox}
You do not ``enable Photon'' on a SQL warehouse---it is the engine. On classic compute you choose a Photon-enabled runtime; on serverless it is always on.
\end{notebox}

\section{Tuesday: Warehouse SQL and the Query Profile}
\subsection{Analytical SQL on the Lakehouse}
Databricks SQL speaks Spark SQL against the three-level namespace, with time travel, window functions, and \texttt{QUALIFY}:

\begin{lstlisting}[language=SQL,caption={Warehouse-grade analytical SQL patterns.},label={lst:ch6-sql}]
-- Top products per region this quarter
SELECT region, product_id, SUM(net_amount) AS revenue
FROM de_training.gold.fact_orders
WHERE order_date >= DATE '2026-07-01'
GROUP BY region, product_id
QUALIFY ROW_NUMBER() OVER (PARTITION BY region ORDER BY revenue DESC) <= 10;

-- Point-in-time question against history (Chapter 2 pays off here)
SELECT COUNT(*) FROM de_training.gold.fact_orders
TIMESTAMP AS OF '2026-08-01T00:00:00Z';
\end{lstlisting}

\subsection{Reading the Query Profile}
Every query has a profile (Query History $\rightarrow$ query $\rightarrow$ profile): a graph of operators with rows, bytes, peak memory, and duration. The diagnostic skill from Chapter 4 transfers directly: look for the operator holding the time---a scan reading more bytes than expected (layout problem, Chapter 5), a spill (memory), or an exchange skewing (join strategy).

\section{Wednesday: Materialized Views, Caching, and Cost Attribution}
\subsection{Materialized Views versus Scheduled Gold}
A DBSQL \textbf{materialized view} is a managed precomputed result with incremental refresh where possible; a \textbf{streaming table} maintains a continuously appended aggregation \cite{databricksSQL}. The decision rule: use a materialized view when consumers query a stable derived shape and freshness tolerance is minutes-to-hours; keep a scheduled Gold table when the derivation is complex business logic that belongs in your pipeline repo with tests (Chapter 3's discipline).

\begin{lstlisting}[language=SQL,caption={A materialized view over the Gold star schema.},label={lst:ch6-mv}]
CREATE MATERIALIZED VIEW de_training.gold.mv_daily_revenue
AS SELECT order_date, region, SUM(net_amount) AS revenue, COUNT(*) AS orders
FROM de_training.gold.fact_orders
GROUP BY order_date, region;
\end{lstlisting}

\subsection{Scaling and Cost}
Concurrency scales by clusters: a 2X-Small warehouse with max 10 clusters absorbs a 9 a.m.\ dashboard burst by adding clusters, then scales back. Auto-stop kills idle spend. Attribution comes from system tables:

\begin{lstlisting}[language=SQL,caption={Warehouse cost attribution from system tables.},label={lst:ch6-cost}]
SELECT w.warehouse_name,
       ROUND(SUM(u.usage_quantity), 2) AS dbus,
       COUNT(DISTINCT u.usage_date) AS active_days
FROM system.billing.usage u
JOIN system.warehouses w USING (warehouse_id)
WHERE u.usage_date >= current_date() - 30
  AND u.sku_name LIKE '%SQL%'
GROUP BY w.warehouse_name
ORDER BY dbus DESC;
\end{lstlisting}

\begin{pitfall}
\textbf{One giant warehouse for everything.} A single maxed-out warehouse mixes the CEO's dashboard with a backfill scan and makes both unpredictable. Split by workload class (executive BI, ad-hoc, engineering), set per-warehouse max clusters, and let each have its own auto-stop.
\end{pitfall}

\section{Thursday Hands-on Project: Serve the Gold Star Schema}
Put the Month 1 Gold schema behind a serverless SQL warehouse, add a materialized view, capture query profiles, and attribute cost.
\index{hands-on lab}\index{SQL warehouse!lab}\index{query profile!lab}

\begin{lstlisting}[language=SQL,caption={Week 6 lab: warehouse setup, MV, and measurement.},label={lst:ch6-lab}]
-- 1. Create warehouse (UI or API): type=serverless, size=2X-Small,
--    min=1, max=4, auto-stop=10 min, channel=Current.

-- 2. Baseline the daily-revenue query directly on the fact table
SELECT order_date, region, SUM(net_amount) AS revenue
FROM de_training.gold.fact_orders
GROUP BY order_date, region;              -- capture the profile

-- 3. Materialized view, then re-run the same question against it
CREATE MATERIALIZED VIEW de_training.gold.mv_daily_revenue AS
SELECT order_date, region, SUM(net_amount) AS revenue, COUNT(*) AS orders
FROM de_training.gold.fact_orders
GROUP BY order_date, region;

SELECT * FROM de_training.gold.mv_daily_revenue
WHERE order_date >= '2026-08-01';         -- capture the faster profile

-- 4. Cost: what did today's work cost?
SELECT sku_name, ROUND(SUM(usage_quantity),2) AS dbus
FROM system.billing.usage
WHERE usage_date = current_date()
GROUP BY sku_name;
\end{lstlisting}

\subsection*{Deliverables}
\begin{itemize}
    \item A running serverless warehouse with documented size, scaling bounds, and auto-stop.
    \item Two captured query profiles (fact scan vs.\ MV) with an explanation of the difference.
    \item The day's DBU attribution query output.
\end{itemize}

\subsection*{Acceptance Criteria}
\begin{itemize}
    \item The MV returns results identical to the direct aggregate (reconciliation).
    \item The warehouse auto-stops within its timer; evidence from the warehouse events tab.
    \item Scaling bounds are justified in one paragraph (expected concurrency, SLA).
    \item The query-profile analysis names the dominant operator in the slow plan.
\end{itemize}

\subsection*{Stretch Goals}
\begin{itemize}
    \item Create a second warehouse for ``engineering'' and prove workload isolation during a simulated burst.
    \item Compare classic vs.\ serverless startup latency with timestamps.
    \item Build the weekly per-warehouse cost report as a saved query with a schedule.
\end{itemize}

\section{Friday Use Case: 500 Analysts, One 9 A.M. Burst, Zero Contention}
A company's 500 analysts open dashboards at 9 a.m.; simultaneously, data engineers run month-end backfills against the same Gold tables. Complaints: dashboards queue, backfills time out, and last month someone ``fixed'' it by doubling warehouse size---and the bill.
\index{capacity planning}\index{workload isolation}

\subsection{Requirements and Assumptions}
Dashboards must load in under 10 seconds at p95; backfills may take hours but must finish by month-end close. Gold tables are liquid-clustered (Chapter 5). Budget grows slower than headcount.

\begin{figure}[H]
    \centering
    \resizebox{\textwidth}{!}{%
    \begin{tikzpicture}[
        wh/.style={rectangle, rounded corners, draw=primary, thick, fill=primary!7, align=center, minimum width=3.4cm, minimum height=1.1cm, font=\small\bfseries},
        load/.style={rectangle, rounded corners, draw=codegray, thick, fill=white, align=center, minimum width=3.0cm, minimum height=0.85cm, font=\scriptsize\bfseries},
        store/.style={cylinder, draw=primary, shape border rotate=90, aspect=0.25, fill=primary!10, align=center, minimum width=2.4cm, minimum height=1.2cm, font=\small\bfseries},
        flow/.style={-{Latex[length=2mm]}, draw=primary, thick}
    ]
        \node[load] (bi) at (0,1.4) {500 analysts\\9 a.m. burst};
        \node[load] (adhoc) at (0,0) {ad-hoc\\analysts};
        \node[load] (eng) at (0,-1.4) {engineering\\backfills};
        \node[wh] (whbi) at (4.6,1.4) {WH-executive\\serverless 2XS\\min 2 / max 10};
        \node[wh] (whadhoc) at (4.6,0) {WH-adhoc\\serverless XS\\min 1 / max 3};
        \node[wh] (wheng) at (4.6,-1.4) {WH-engineering\\pro Small\\min 1 / max 2};
        \node[store] (gold) at (9.6,0) {Gold Delta\\(liquid clustered)};
        \draw[flow] (bi) -- (whbi);
        \draw[flow] (adhoc) -- (whadhoc);
        \draw[flow] (eng) -- (wheng);
        \draw[flow] (whbi) -- (gold);
        \draw[flow] (whadhoc) -- (gold);
        \draw[flow] (wheng) -- (gold);
    \end{tikzpicture}%
    }
    \caption{Workload-class warehouse topology over one governed Gold layer: isolation without data copies.}
    \label{fig:ch6-warehouse-topology}
\end{figure}

\subsection{Architecture Decisions}
\textbf{Three warehouses, three classes.} Executive BI (serverless, aggressive scaling, tight auto-stop), ad-hoc (small, cheap, forgiving), engineering (pro, sized for scans, auto-stop generous). The shared Delta storage means isolation costs no copies; the burst is absorbed by max-clusters scaling, not by an oversized idle floor.

\textbf{Query hygiene.} Dashboards query the materialized view, not the raw fact; the backfill writes at off-peak with its own warehouse; a saved monitoring query (profiles + billing) reviews p95 weekly. The doubling-the-size incident becomes a scaling-bounds conversation backed by queue-time metrics.

\begin{pitfall}
\textbf{Reading the queue as a sizing problem.} Queued queries at 9 a.m.\ with idle clusters mean the \emph{max clusters} bound is the constraint, not size; a bigger single cluster would not have helped and would have billed more. Queue time, running time, and cluster count are three different metrics---read all three before changing anything.
\end{pitfall}

\subsection{Risks and Mitigations}
\begin{table}[H]
\centering
\small
\begin{tabular}{@{}p{4.6cm}p{7.6cm}@{}}
\toprule
\textbf{Risk} & \textbf{Mitigation} \\
\midrule
Burst budget spikes with headcount & Max-clusters caps per warehouse; monthly chargeback review \\
A heavy query starves the executive warehouse & Workload classes isolated; per-user query limits \\
Stale MV presented as fresh & Staleness alert; refresh timestamp rendered on the dashboard \\
Warehouse sprawl across teams & Inventory with owners; unused warehouses auto-flagged \\
Classic warehouses forgotten running & Auto-stop enforced by policy; weekly idle audit \\
\bottomrule
\end{tabular}
\caption{Week 6 scenario risks and mitigations.}
\label{tab:ch6-risks}
\end{table}

\section{Design Decision Guide}
\index{decision guide!Week 6}%
Warehouse decisions are workload-shape decisions; name the workload class first.

\begin{table}[H]
\centering
\small
\begin{tabular}{@{}p{3.4cm}p{4.4cm}p{4.6cm}@{}}
\toprule
\textbf{Decision} & \textbf{Choose the first when} & \textbf{Choose the second when} \\
\midrule
Serverless vs.\ classic warehouse & BI and bursty SQL; zero-ops valued & Custom networking or lowest-rate steady load \\
Materialized view vs.\ scheduled Gold & Stable derived shape, tolerance for managed refresh & Complex tested business logic \\
One warehouse vs.\ per-class warehouses & Tiny teams, light concurrency & Mixed latency classes or chargeback needs \\
Size up vs.\ scale out & Individual queries are slow & Concurrency (queueing) is the problem \\
Result cache reliance vs.\ MV & Ad-hoc repeated questions & Shared, governed derived numbers \\
Photon confidence vs.\ measurement & Always measure & --- \\
\bottomrule
\end{tabular}
\caption{Week 6 decision guide.}
\label{tab:ch6-decisions}
\end{table}

The recurring trap is treating warehouse size as the only knob. Bounds, classes, caching, and derived tables usually matter more---and all four are cheaper than a bigger T-shirt size.

\section{Operational Guidance for Week 6}
\index{operational guidance!Week 6}%
\subsection{Performance and Reliability}
Review warehouse telemetry weekly: queue time (concurrency bound), median/p95 run time (query health), and cluster count over the day (scaling shape). Auto-stop settings are audited against reality---a warehouse ``stopped at night'' that shows 2 a.m.\ activity is running something nobody scheduled. Materialized-view refresh status is monitored like any pipeline; a silently stale MV is a certified number quietly rotting.

\subsection{Security and Governance Checklist}
\begin{itemize}
    \item Dashboards and scheduled queries publish under a least-privilege service identity.
    \item Warehouse access granted per workload class; engineers do not run backfills on the executive warehouse.
    \item Every warehouse has a documented owner, size rationale, and auto-stop value.
    \item Monthly cost attribution per warehouse is reviewed and tagged to a domain.
\end{itemize}

\subsection{Troubleshooting Playbook}
\begin{table}[H]
\centering
\small
\begin{tabular}{@{}p{3.6cm}p{4.2cm}p{4.8cm}@{}}
\toprule
\textbf{Symptom} & \textbf{Likely cause} & \textbf{First action} \\
\midrule
Queries queue at peak & Max-clusters bound reached & Read queue vs.\ run time; raise max for the burst window \\
First dashboard load slow & Cold warehouse (auto-stop) & Serverless type or a scheduled early-morning warm query \\
MV returns stale rows & Refresh not scheduled/failed & Check refresh history; alert on staleness \\
Same query fast then slow & Result cache invalidated by new commits & Expected; cache hits only when data and text are unchanged \\
\bottomrule
\end{tabular}
\caption{Week 6 troubleshooting quick reference.}
\label{tab:ch6-trouble}
\end{table}

\section{Interview Q\&A}
\subsection{Concepts and Trade-offs}
\begin{enumerate}
    \item \textbf{Q: Why is Databricks SQL different from connecting a BI tool to Spark?}\\
    A: Photon-vectorized execution, warehouse-level caching and scaling, governance integration, and per-query profiles---an analytical serving layer, not a general cluster with a SQL cell.
    \item \textbf{Q: Serverless warehouse: what do you give up?}\\
    A: Data-plane networking options and some custom init control; you gain instant start, fast autoscaling, and no cluster management. For BI, the trade is almost always right.
    \item \textbf{Q: Materialized view or a scheduled Gold table?}\\
    A: MV for stable derived shapes with freshness tolerance and simple logic; scheduled Gold when the derivation is versioned business logic that needs tests, review, and pipeline discipline.
    \item \textbf{Q: How do you find who spent what on SQL last month?}\\
    A: \texttt{system.billing.usage} joined to \texttt{system.warehouses}, broken out by SKU and day; add \texttt{system.access.audit} for per-user query attribution.
    \item \textbf{Q: A dashboard is slow only at 9 a.m. Diagnose.}\\
    A: Queue time at the concurrency bound: raise max clusters for the burst window, move the dashboard to the MV, and check that the warehouse is serverless so scale-up is seconds, not minutes.
    \item \textbf{Q: When is a classic warehouse still the right answer?}\\
    A: Custom networking/init requirements, the strictest per-DBU price on steady loads, or features gated to the classic tier---decided from requirements, not habit.
    \item \textbf{Q: What does the query profile tell you that the UI duration does not?}\\
    A: \emph{Where} the time went: which operator scanned how many bytes and spilled what. Duration says slow; the profile says why.
    \item \textbf{Q: How do you prove a warehouse right-sizing saved money?}\\
    A: Same workload, measured weeks, \texttt{system.billing.usage} deltas per SKU---the cost chapter of the tuning report, with queue time held constant so the saving is not just slower service.
\end{enumerate}

\section{Further Reading}
The Databricks SQL \cite{databricksSQL} and SQL warehouse \cite{databricksSQLWarehouses} documentation cover types, sizing, and monitoring; Photon is explained in \cite{databricksPhoton}. System-table billing and warehouse monitoring are documented in \cite{databricksSystemTables}. For the engine fundamentals, revisit Spark SQL tuning \cite{sparkSQLTuning}.

\section*{Key Takeaways}
\begin{itemize}
    \item A SQL warehouse is a scalable query fleet over governed Delta tables---no copies, no separate store.
    \item Serverless is the BI default; pro/classic serve networking and control edge cases.
    \item Materialized views answer repeated derived questions; complex business logic still belongs in tested pipelines.
    \item Queue time, run time, and cluster count are different metrics---size nothing before reading all three.
    \item Cost attribution is a system-table query away; make it a weekly habit.
\end{itemize}

\section{Exercises}
\begin{enumerate}
    \item \textbf{(Conceptual)} Contrast a SQL warehouse with an all-purpose cluster for a BI workload: three differences.
    \item \textbf{(Conceptual)} Why does the result cache make some dashboard refreshes free, and when does it miss?
    \item \textbf{(Hands-on)} Complete the Thursday lab, including both query profiles.
    \item \textbf{(Hands-on)} Build the three-warehouse topology from the Friday case and measure a simulated burst (20 concurrent queries) against each.
    \item \textbf{(Hands-on)} Write and schedule the weekly warehouse cost report; include a week-over-week delta column.
    \item \textbf{(Design)} Choose warehouse type/size/bounds for: executive dashboards, a data-science team, and a dbt-style transformation workload. Justify each.
    \item \textbf{(Operations)} Write the alert rule for ``warehouse queue time $>$ 60 s for 10 minutes.''
    \item \textbf{(Architecture)} When would you put a materialized view in a DLT pipeline instead of DBSQL (preview of Chapter 10)?
\end{enumerate}
